\section{Pilar A: detección de secuencias multi-etapa}
\label{sec:pilarA}

La detección basada en reglas descrita en la Sección~\ref{sec:deteccion-reglas} produce
alertas aisladas: cada una señala una técnica concreta en un instante concreto. El Pilar~A
transforma esas alertas en \textbf{campañas}, es decir, en conjuntos de evidencias que el
sistema reconoce como etapas sucesivas de un mismo ataque. Esta sección describe el corpus
sobre el que se ha construido, el método de validación, la estructura de correlación en dos
niveles que impone la herramienta, los tres patrones implementados, el esquema común en que
se emiten las campañas y su validación en modo continuo.

\subsection{Corpus de evaluación y verdad conocida}

Todo el trabajo de esta fase se apoya en un \textbf{corpus fijo y reproducible}, cargado
mediante un único guion que recrea los \textit{topics}, reproduce los conjuntos de datos y
los procesa con grupos de consumo nuevos (Anexo~\ref{anx:reproducibilidad}). El corpus reúne
nueve conjuntos de datos: el primer día de la emulación de APT29 publicada por MITRE, una
captura atómica de volcado de LSASS mediante \texttt{comsvcs.dll} y siete campañas de robo
de credenciales que emplean siete herramientas distintas.

\begin{table}[H]
\centering
\begin{tabular}{|l|r|}
\hline
\textbf{Magnitud} & \textbf{Valor} \\
\hline
Eventos consumidos & 550.494 \\
Eventos normalizados a OCSF & 530.401 (96,4\,\%) \\
Eventos derivados a la cola de rechazos & 20.093 \\
Reglas Sigma evaluadas & 2.093 \\
Alertas distintas producidas & 1.184 \\
\quad de ellas, en el primer día de APT29 & 741 \\
\hline
\end{tabular}
\caption{Corpus de evaluación del Pilar A.}
\label{tab:corpus-f3}
\end{table}

La totalidad de los rechazos corresponde a tipos de evento sin mapeador, y en torno al 60\,\%
son marcadores de inicio y fin de bloques de PowerShell sin contenido. Ningún evento de
acceso a procesos ni de creación de ficheros figura entre ellos, de modo que la evidencia
central del robo de credenciales se normaliza sin pérdida.

La \textbf{verdad conocida} del primer día de APT29 se reconstruyó a partir del plan de
emulación publicado: 28 pasos agrupados en 10 etapas, con su equipo y usuario. Los
identificadores de técnica del plan corresponden a versiones antiguas de ATT\&CK y varios
están obsoletos (por ejemplo, \texttt{T1086} es hoy \texttt{T1059.001} y \texttt{T1088} es
\texttt{T1548.002}). Sin traducirlos, el cruce entre el plan y las alertas no cuadra y
aparenta fallos de detección que no lo son.

\subsection{Método: versión de referencia y versión en flujo}

Cada patrón se ha implementado \textbf{dos veces}. La primera es una vista SQL sobre el
almacén analítico que procesa el corpus completo por lotes y actúa como
\textbf{oráculo}: fija el resultado esperado. La segunda es la implementación en Flink con
\texttt{MATCH\_RECOGNIZE}, que debe reproducir ese resultado.

El valor de este método es separar las fuentes de error. Sin oráculo, depurar el patrón
obliga a dudar a la vez de su lógica y de la semántica temporal del motor; con él, la
pregunta se reduce a por qué no aparece lo que se sabe que existe. El método funciona en
ambos sentidos: en un caso fue la versión en flujo la que reveló un defecto del oráculo, que
descartaba la captura atómica por exigir contexto alrededor de las evidencias, requisito que
penaliza precisamente las capturas más limpias.

Para validar en Flink sin depender del avance de las marcas de agua, la fuente del bus se
configura como \textbf{acotada} (se lee hasta el último desplazamiento disponible) y el
motor se ejecuta en modo por lotes. Así cada consulta produce un resultado final, comparable
fila a fila con el oráculo. La validación en modo continuo se aborda por separado en la
Sección~\ref{sec:pilarA-streaming}.

\subsection{Correlación en dos niveles}
\label{sec:dos-niveles}

La intuición inicial era expresar cada campaña con un único patrón: «todas las evidencias
consecutivas de un mismo equipo hasta que se abra un hueco de más de diez minutos». Esa
expresión resulta \textbf{imposible} con \texttt{MATCH\_RECOGNIZE} en Flink. El motor no
admite cuantificadores voraces como último elemento de un patrón, y el cuantificador
reluctante, que sí admite, produce la coincidencia \textit{mínima}: con
\texttt{PATTERN (E+? C)} se obtienen 29 coincidencias de dos filas donde el oráculo
reconoce 10 campañas.

La solución es una \textbf{correlación jerárquica en dos niveles}, que aquí deja de ser una
opción de diseño para convertirse en un requisito impuesto por la herramienta:

\begin{enumerate}
    \item \textbf{Nivel 1, transiciones.} \texttt{MATCH\_RECOGNIZE} reconoce pares de
    evidencias consecutivas que cumplen la condición del patrón dentro de una ventana
    temporal. Es lo que un motor CEP hace bien.
    \item \textbf{Nivel 2, campañas.} Una agregación agrupa las transiciones separadas por
    un hueco menor que la ventana. Es lo que un motor CEP hace mal y una agregación resuelve
    de forma trivial. En lotes se implementa con funciones de ventana analíticas; en flujo,
    con una ventana de sesión sobre el tiempo de cada transición.
\end{enumerate}

\begin{figure}[H]
\centering
\begin{tikzpicture}[
  ev/.style={circle,draw=black!70,fill=black!10,minimum size=4mm,inner sep=0pt},
  evc/.style={ev,fill=black!55},
  tr/.style={-{Stealth[length=2mm]},thick},
  font=\footnotesize]
  \foreach \x/\t in {0/ev,1.1/evc,2.0/ev,3.2/evc,7.4/ev,8.4/evc,9.3/ev} {
    \node[\t] at (\x,0) {};
  }
  \draw[tr] (0.15,0.25) to[bend left=40] (0.95,0.25);
  \draw[tr] (1.25,0.25) to[bend left=40] (1.85,0.25);
  \draw[tr] (2.15,0.25) to[bend left=40] (3.05,0.25);
  \draw[tr] (7.55,0.25) to[bend left=40] (8.25,0.25);
  \draw[tr] (8.55,0.25) to[bend left=40] (9.15,0.25);
  \node[anchor=west] at (9.8,0.45) {Nivel 1: transiciones};
  \draw[thick] (-0.2,-0.45) -- (3.4,-0.45);
  \draw[thick] (-0.2,-0.35) -- (-0.2,-0.55); \draw[thick] (3.4,-0.35) -- (3.4,-0.55);
  \draw[thick] (7.2,-0.45) -- (9.5,-0.45);
  \draw[thick] (7.2,-0.35) -- (7.2,-0.55); \draw[thick] (9.5,-0.35) -- (9.5,-0.55);
  \node at (1.6,-0.8) {campaña 1};
  \node at (8.35,-0.8) {campaña 2};
  \draw[<->,dashed] (3.4,-0.15) -- node[above] {hueco $>$ ventana} (7.2,-0.15);
  \node[anchor=west] at (9.8,-0.45) {Nivel 2: agregación por hueco};
  \draw[->] (-0.4,-1.25) -- (12.8,-1.25) node[right] {$t$};
\end{tikzpicture}
\caption{Correlación en dos niveles. Los círculos oscuros son evidencias de la etapa ancla
del patrón; las flechas, las transiciones reconocidas por \texttt{MATCH\_RECOGNIZE}.}
\label{fig:dos-niveles}
\end{figure}

La implementación encontró además otras restricciones del motor, que condicionan la forma
de escribir los patrones: la cláusula \texttt{DEFINE} no admite subconsultas, por lo que la
lógica sobre listas de técnicas se precalcula en la proyección; \texttt{COUNT(DISTINCT)} no
está permitido en \texttt{MEASURES}; y \texttt{WITHIN} acota la duración total de la
coincidencia y no el hueco entre elementos, por lo que la ventana se expresa como condición
en \texttt{DEFINE}.

Una restricción adicional determina la arquitectura. \texttt{MATCH\_RECOGNIZE} exige ordenar
por un atributo de tiempo de evento, y Flink lo pierde al atravesar uniones y agregaciones.
Por ello, cada patrón se divide en dos trabajos desacoplados por el bus: uno de
\textbf{preparación} publica en un \textit{topic} intermedio la evidencia del patrón, sin
repeticiones y en orden de tiempo de evento, y otro de \textbf{reconocimiento} la relee con
su marca de agua. Es también la arquitectura adecuada para producción, pues separa el
enriquecimiento del reconocimiento. El Anexo~\ref{anx:patrones} recoge el código del
patrón P1 como ejemplo.

\subsection{P1: robo de credenciales}

P1 reconoce la secuencia herramienta $\rightarrow$ acceso a la memoria de LSASS
$\rightarrow$ volcado $\rightarrow$ limpieza. Su diseño se decidió con datos en tres puntos.

\textbf{El ancla es la familia \texttt{T1003.*} y no la subtécnica \texttt{T1003.001}.}
Medido sobre las siete variantes, la subtécnica exacta cubre seis: falla precisamente
Mimikatz, la variante más emblemática, porque Meterpreter lo ejecuta dentro de su propio
proceso y ningún proceso externo abre la memoria de LSASS. La familia cubre las siete.

\textbf{Las etapas se reconocen por el tipo de evento de origen}, restringido a alertas de
la familia \texttt{T1003}: un acceso a proceso es acceso a LSASS, una creación de proceso es
ejecución de la herramienta y una creación de fichero es volcado a disco. Una primera versión
reconocía las etapas por los campos coincidentes de la regla y producía evidencia falsa: las
292 alertas de un \textit{dropper} de PowerShell de APT29 satisfacían la etapa «volcado a
fichero» sin relación alguna con credenciales. Con la restricción, la evidencia de P1 en el
equipo SCRANTON pasa de 308 alertas a 3.

\textbf{La transición es bidireccional}: basta con que una de sus dos evidencias sea de
credenciales, de modo que valen tanto «preparación $\rightarrow$ credenciales» como
«credenciales $\rightarrow$ limpieza». Con la versión unidireccional se perdía el equipo
NASHUA, donde la limpieza llega después del único acceso a credenciales. La ventana es de
10 minutos, el doble del núcleo observado en APT29 (4\,min 47\,s entre la descarga de la
herramienta y el acceso a LSASS).

La confianza se \textbf{descompone en factores}: 0,40 por el ancla, más 0,25 si hay acceso a
LSASS, 0,20 si hay volcado a disco, 0,10 si hay ejecución de la herramienta y 0,15 si hay
limpieza, con un máximo de 1. El resultado son \textbf{11 campañas}: una por cada una de
las \textbf{siete variantes} de robo de credenciales, dos en SCRANTON y una en NASHUA dentro
de APT29, y la captura atómica. La confianza oscila entre 0,50 (Mimikatz, solo ancla y
contexto) y 0,95 (las variantes con acceso a LSASS y fichero de volcado).

\subsection{P2: movimiento lateral}

P2 reconoce el salto de un equipo a otro: una conexión de red hacia el destino seguida de
ejecución remota en él. Es el primer patrón que \textbf{cruza equipos}, y requiere
construir antes las \textbf{aristas} de red a partir de los eventos normalizados.

La construcción de las aristas destapó varios problemas que no producían error alguno.
Bajo la misma clase OCSF de actividad de red conviven las conexiones lógicas de Sysmon y
cada flujo permitido por el cortafuegos de Windows; filtrar por clase daba 3.480 conexiones
WinRM donde el fichero de origen tiene 12. El mapa de direcciones IP a equipos, si cuenta
también los registros entrantes, se corrompe en silencio: un equipo llegó a ganar su propia
dirección por un solo voto. Se resolvió contando solo conexiones salientes, en las que la
dirección de origen es por definición la propia. Y la dirección de bucle local escrita sin
abreviar fabricaba aristas falsas hasta que se normalizó su grafía.

Con el mapa corregido y el controlador de dominio identificado por su comportamiento y
excluido, quedan 39 aristas: 22 hacia el controlador y \textbf{17 de SCRANTON a NASHUA}
(puertos 135, 445 y 5985), todas confirmadas desde los dos extremos. Sobre ellas, el nivel~1
reconoce en cada destino una arista y una evidencia de ejecución (técnicas \texttt{T1021},
\texttt{T1569.002} y \texttt{T1003}) separadas como máximo 300 segundos, en cualquier orden,
y el nivel~2 agrupa las transiciones de un mismo par de equipos.

El resultado es \textbf{una campaña}, idéntica campo a campo a la del oráculo: de las
03:09:23.801 a las 03:17:59.493, con 8 evidencias en el destino, cuatro etapas (conexión,
ejecución remota, credenciales y servicio) y confianza 1,0. El orden no se exige pero se
\textbf{mide}: siete de las ocho evidencias llegan después de su conexión, con un retraso
máximo de 1.004\,ms. La octava tiene una hora de resolución de segundo y no permite ordenar
dentro de ese segundo.

\subsection{P3: escalada de privilegios con descubrimiento o persistencia}

El inventario previo mostró que \textbf{las reglas no ven el descubrimiento del atacante}
en APT29: solo hay dos alertas de la táctica en todo el día, y ambas proceden del agente de
la nube del equipo. El descubrimiento real está en los bloques de PowerShell registrados
(\texttt{Invoke-Discovery}, \texttt{Get-Process} remoto). Por ello P3 es el primer patrón
\textbf{híbrido}: su etapa de descubrimiento combina alertas Sigma y eventos de PowerShell
con instrucciones de reconocimiento.

El ancla es la \textbf{escalada de privilegios} (técnicas \texttt{T1548}, \texttt{T1134},
\texttt{T1055} y \texttt{T1068}), acompañada de descubrimiento o de persistencia
(\texttt{T1547}, \texttt{T1543}, \texttt{T1053}, \texttt{T1546}, \texttt{T1037} y
\texttt{T1136}) en el mismo equipo. Los eventos de asignación de privilegios especiales se
descartaron como etapa porque aparecen en todos los equipos de forma rutinaria. Se aplica
una única exclusión, medida: la modificación de claves de arranque que escribe el
explorador de Windows al iniciar sesión, que eliminaba 28 evidencias espurias. Una segunda
exclusión prevista para los agentes de la nube se retiró al comprobar que no eliminaba
nada: una exclusión sin efecto es un ajuste especulativo.

El resultado son \textbf{3 campañas}. En APT29, la de SCRANTON (confianza 0,90) reproduce
el orden exacto del plan del atacante: escalada $\rightarrow$ descubrimiento $\rightarrow$
persistencia. Las otras dos son las capturas de robo de credenciales que incluyen escalada
real (confianza 0,75). Lo que el patrón \textbf{deja fuera} es tan relevante como lo que
detecta: el controlador de dominio, con 25 evidencias de descubrimiento producidas por un
guion de monitorización legítimo, y cinco capturas con entre 26 y 64 evidencias de
reconocimiento no forman campaña porque no hay escalada. El software legítimo hace el mismo
reconocimiento que un atacante; es la secuencia lo que los distingue.

\begin{table}[H]
\centering
\small
\begin{tabular}{|l|p{4.3cm}|p{5.2cm}|}
\hline
\textbf{Patrón} & \textbf{Qué reconoce} & \textbf{Resultado} \\
\hline
P1 Robo de credenciales & Herramienta, acceso a LSASS, volcado y limpieza & 11 campañas; las 7 variantes de robo de credenciales detectadas \\
\hline
P2 Movimiento lateral & Conexión a otro equipo y ejecución remota en el destino & 1 campaña SCRANTON $\rightarrow$ NASHUA, 17 aristas bilaterales \\
\hline
P3 Escalada & Escalada con descubrimiento o persistencia & 3 campañas; orden del plan de APT29 reproducido \\
\hline
\end{tabular}
\caption{Patrones del Pilar A y resultados sobre el corpus. En los tres, la versión en
Flink reproduce exactamente su oráculo.}
\label{tab:patrones}
\end{table}

\subsection{Esquema común de campaña}

Las campañas de los tres patrones se emiten con un \textbf{esquema común}, que es lo que
consumirán el grafo de ataque, la priorización y la interfaz. Por coherencia con el resto
del sistema, una campaña es un hallazgo estándar \textbf{OCSF 1.5.0 «Detection Finding»}
(clase 2004). Todo lo que el estándar contempla va en sus campos: las alertas relacionadas,
las técnicas, la confianza y el intervalo temporal. Lo que no contempla va en una extensión
propia documentada: las \textbf{etapas ordenadas}, los \textbf{factores de confianza}, cuya
suma es la confianza publicada, los equipos con su papel y las métricas de cada patrón. Las
técnicas se publican sin táctica asociada, porque las alertas traen ambas listas sin
emparejar y emparejarlas sería inventar.

Dos reglas garantizan que el resultado sea reproducible. La identidad de cada campaña es un
resumen criptográfico de su conjunto de datos, patrón, clave y hora de inicio, de modo que
reprocesar el corpus produce los mismos identificadores y la tabla de destino sobrescribe en
lugar de duplicar. Las etapas se ordenan por el segundo de su primera evidencia y, dentro
del mismo segundo, por el orden lógico del patrón, porque parte de las horas de origen solo
tienen resolución de segundo. El proceso que vuelca las campañas al almacén
\textbf{valida cada documento} y se detiene ante el primero que no cumpla el esquema. El
almacén contiene hoy 15 campañas: 11 de P1, 1 de P2 y 3 de P3. El
Anexo~\ref{anx:esquema} detalla el esquema con un ejemplo completo.

\subsection{Validación en modo continuo}
\label{sec:pilarA-streaming}

La validación anterior se realizó con fuentes acotadas. En modo continuo, la marca de agua
descarta sin aviso los eventos que llegan con más retraso del admitido, por lo que antes de
lanzar los patrones se \textbf{midió el desorden} de llegada:

\begin{itemize}
    \item Dentro de un mismo conjunto de datos, las alertas llegan con hasta 293\,s de
    retraso respecto a otras posteriores, por efecto de los detectores en paralelo. Con una
    marca de agua de 30\,s se perderían 220 alertas (18\,\%); con 5 minutos, ninguna. La
    cifra está inflada porque el corpus se reproduce decenas de veces más rápido que el
    tiempo real.
    \item Los \textit{topics} intermedios, escritos inicialmente sin ordenar, habrían
    llegado «tarde» en un 74--78\,\%. Se decidió publicarlos en orden de tiempo de evento.
    \item Al reproducir a la vez capturas de distintos años, el desorden entre conjuntos de
    datos alcanza años. Es un artefacto de la reproducción que no existe en producción.
\end{itemize}

Con los \textit{topics} intermedios ordenados, el segundo nivel se expresa con una ventana de
sesión y los tres patrones dan en flujo \textbf{exactamente el mismo resultado} que por lotes:
las 11 campañas de P1 con sus 52 transiciones, la campaña de P2 y las 3 de P3. La ventana de
sesión no emite cada campaña una sola vez al cerrarse, sino que la \textbf{actualiza} con cada
transición nueva. Para detectar es una ventaja, porque la campaña existe desde su primera
transición, pero obliga a emitirla hacia un destino que admita actualizaciones.

\subsection{Incidencia: la hora de los eventos era la de recogida}

El primer día de APT29 no incluye el campo de hora de creación del evento, y la
normalización recurría a la marca del colector, que corresponde al instante de
\textbf{recepción}. El desfase medido sobre los eventos que sí traen la hora real tiene una
mediana de 6,2\,s, un percentil 95 de 74,9\,s y un máximo de 101,2\,s. Lo grave es su
dispersión, de 26,7\,s: los eventos se reordenaban entre sí y no había desplazamiento
constante que corregir.

La corrección establece una cadena de prioridad entre las fuentes de tiempo disponibles, con
la conversión de zona horaria derivada por conjunto de datos, y marca cada evento con la
fuente empleada. Ninguna fuente es fiable en todos los conjuntos: la hora UTC de Sysmon, la
más precisa en APT29, va desviada 19 horas en las capturas de robo de credenciales. Con la
hora corregida desapareció una inversión imposible documentada hasta entonces, en la que el
proceso remoto aparecía tres segundos antes que la conexión que lo provoca: la conexión va
ahora 116\,ms antes. Queda una limitación del dato de origen: el 26,6\,\% de los eventos de
APT29 solo tiene resolución de segundo.

\subsection{Incidencia: la identidad de los eventos y de las alertas}

Durante esta fase se comprobó que el 15\,\% de los eventos normalizados de APT29 era idéntico
byte a byte a otro, porque el evento OCSF no conservaba el número de registro de origen. Dos
ocurrencias distintas con los mismos campos se volvían indistinguibles, y las alertas
duplicadas por un reprocesado no se podían separar de las legítimas. Se resolvió dando
\textbf{identidad determinista} a ambos: el evento conserva su número de registro (o, si la
fuente no lo aporta, el ordinal de la línea en el fichero), la alerta se identifica por el
resumen del evento y la regla, y las tablas del almacén deduplican por esa identidad. La
comprobación final fue recuperar exactamente las 741 alertas documentadas para APT29 y las
cifras de eventos de cada conjunto de datos contadas contra su fichero de origen.

\subsection{Estado de la fase y trabajo pendiente}

Se cumplen los criterios de cierre relativos a la detección: tres patrones que disparan sobre
datos reales, las siete variantes de robo de credenciales detectadas y campañas multi-etapa
reconocidas en APT29, con resultados idénticos por lotes y en flujo. Quedan pendientes:

\begin{itemize}
    \item Medir los \textbf{falsos positivos} sobre la porción benigna del corpus y sobre
    capturas atómicas de una sola táctica.
    \item Medir la \textbf{ratio de reducción de alertas} de los tres patrones. La de P1
    (0,991) no es comparable con la línea base de agrupación de la fase anterior, porque solo
    cubre lo que P1 absorbe alrededor de sus anclas.
    \item Decidir la \textbf{extensión de una campaña}: la versión en Flink la acota a su
    evidencia y el oráculo absorbe $\pm$10 minutos. En P3 la diferencia es grande, porque el
    descubrimiento llega en ráfagas alrededor del ancla, y condiciona la ratio de reducción.
    \item Resolver el \textbf{punto ciego del etiquetado}: el 27,1\,\% de las alertas no lleva
    técnica ATT\&CK y es invisible para la correlación por construcción.
    \item Incorporar el segundo día de APT29 y emitir las campañas en flujo hacia un destino
    que admita actualizaciones.
\end{itemize}
